%% ============================================================
%% Simulation Methodology Section
%% ============================================================

\section{Simulation Methodology}
\label{sec:simulation}

\subsection{Network Topology}

We simulate a heterogeneous NTN graph $G=(V,E)$ with
$N=1{,}000$ nodes across five segment types: GEO satellites
($n_{\mathrm{GEO}}=2$), MEO satellites ($n_{\mathrm{MEO}}=4$),
LEO satellites ($n_{\mathrm{LEO}}=30$), HAPS nodes
($n_{\mathrm{HAPS}}=10$), UAVs ($n_{\mathrm{UAV}}=50$), and
IoT ground devices ($n_{\mathrm{Ground}}=904$).  Links are generated
from a segment-aware random graph model with per-link propagation
delays drawn uniformly from segment-typical ranges
(GEO: 150--300\,ms; LEO: 5--40\,ms; Ground: 1--20\,ms).  Connectivity
is guaranteed via a minimum spanning-tree augmentation step.
All topology parameters are fixed in \texttt{configs/simulation\_config.yaml}
to ensure byte-for-byte reproducibility.

\subsection{Trust Update and Baselines}

Each simulation consists of $R=500$ rounds of duration $T_{\mathrm{round}}=60$\,s,
representing one NTN interaction cycle (aligned with 3GPP NTN
scheduling periodicity for IoT devices~\cite{3gpp_ntn}).
In each round, every node $i$ observes one randomly selected neighbour $j$
and computes a feedback score $s_{ij}(t) \in [0,1]$ based on
packet-delivery ratio and authentication success.  The feedback is
submitted to the DLT, which applies Eq.~(4):
%
\begin{equation*}
  T_{ij}(t+\Delta t) =
    \alpha \cdot T_{ij}(t) \cdot e^{-\lambda_s \Delta t}
    + (1-\alpha)\,s_{ij}(t)
\end{equation*}
%
where $\Delta t = T_{\mathrm{round}} + d_{ij}/1000$ (seconds), $d_{ij}$
is the link propagation delay in milliseconds, and $\lambda_s$ is the
segment-specific decay rate (Table~\ref{tab:lambda}).

Five trust models are compared:
\begin{enumerate}
  \item \textbf{Proposed}: Eq.~(4) with per-segment $\lambda_s$,
        Paillier HE aggregation, PBFT-like DLT, and ZTA policy layer.
  \item \textbf{Centralized Trust}: Authority aggregates all feedback;
        inaccessible during partition (simulates single-authority SPOF).
  \item \textbf{Static DLT (NxGenT)}~\cite{nxgent}: Permissioned DLT with
        monotonic trust accumulation; no temporal decay.
  \item \textbf{ZT Auth-Only}: Binary ZTA authentication without a
        continuous trust score; node classified by latest authentication
        outcome.
  \item \textbf{UAV+FL Blockchain}~\cite{uav_fl}: Federated-learning-based
        trust aggregation with uniform $\lambda_{\mathrm{UAV}}=0.05$ decay
        applied to all segment types.
\end{enumerate}

\subsection{Attack Scenarios}

Twenty percent of nodes (unless varied in Exp6) are designated
malicious with roles drawn uniformly from:
\textit{on-off} (alternating good/bad behaviour every 50 rounds),
\textit{build-then-betray} (legitimate for rounds 1--150, malicious
thereafter),
\textit{Sybil} (colluding group with shared false feedback),
and \textit{badmouthing} (transmits falsely low scores for legitimate
peers).  Feedback scores for malicious nodes are:
$s_{ij}(t) = \min(0.15, \mathcal{N}(0.1, 0.05^2))$ during attack
phases and $s_{ij}(t) \sim \mathcal{U}(0.7, 1.0)$ during good
phases.

\subsection{Evaluation Metrics}

\begin{itemize}
  \item \textbf{Detection Latency}: Rounds elapsed from attack onset
        until the malicious node's trust score falls continuously below
        $\theta$ for three consecutive rounds (averaged over all malicious
        nodes and all 30 runs).
  \item \textbf{F1 Score / AUC-ROC}: Binary classification metrics at
        the final round ($t=500$) using threshold $\theta$ and the full
        ROC curve respectively.
  \item \textbf{Trust MAE}: Mean absolute error between estimated and
        ground-truth trust scores (used for privacy-utility tradeoff).
  \item \textbf{Ledger TPS}: Simulated transaction throughput of the
        PBFT-like consensus layer.
\end{itemize}

\subsection{Reproducibility}

All experiments are seeded with \texttt{seed=42} and controlled
via \texttt{configs/simulation\_config.yaml}.  The full pipeline is
executed as:
%
\begin{verbatim}
  pip install -r requirements.txt
  python run_all.py
\end{verbatim}
%
A GitHub Actions CI workflow automatically re-executes the 16-test
unit suite and validates key invariants (trust boundedness,
monotone decay ordering, Paillier correctness) on every commit.
Numerical results are archived as JSON files in
\texttt{results/data/}.  A Zenodo DOI is provided for
long-term archiving of the exact artefact set~\cite{zenodo2026}.
